\section{Introduction}

EEG foundation models aim to provide features that can be reused across
recording systems, electrode layouts, and prediction tasks. REVE addresses
this goal through large-scale pretraining and position encodings that support
different electrode arrangements \citep{elouahidi2025reve}. Evaluating a frozen
encoder, however, also requires choosing a prediction head. A linear head tests
whether the target can be recovered by a simple readout; a more expressive head
may exploit additional structure in the features, but its gains may depend on
the training data.

Age prediction offers a useful test. Age-related changes are reflected in EEG,
and previous work has used them to study maturation and aging
\citep{sun2019brainage,vandenbosch2019age}. Yet a model trained on one cohort may
perform differently on another because participants, hardware, and recording
procedures change. Repeating training with several random seeds measures one
source of variability, while all runs may still be evaluated on the same small
group of people.

We ask whether richer heads improve external age prediction when the REVE
encoder, development data, preprocessing, training budget, and checkpoint
selection rule are held fixed. We compare a mean-pooled linear baseline with a
linear probe of the preceding layer and two richer final-layer heads. All heads
are developed on the Healthy Brain Network (HBN) and evaluated on a held-out
MIPDB cohort using a decision rule specified before external evaluation.

Three secondary analyses examine training-set size, representation depth, and
transfer to a separate cohort, OpenNeuro ds006780. The first two reuse MIPDB;
the third adds new participants but was motivated by the earlier results.

The contribution is a controlled empirical study of frozen EEG features. We
show how the preferred head changes with the metric, how a candidate's
advantage can change with training-set size, and why agreement across seeds
can coexist with substantial uncertainty across participants. These findings
help interpret probe comparisons without attributing every downstream gain
to the encoder itself.
